常见数论函数包括常数函数
\[
  \mathbf 1(n)=1,
\]
恒等函数
\[
  \operatorname{id}(n)=n,
\]
单位函数
\[
  \varepsilon(n)=\begin{cases}1,&n=1,\\0,&n>1,\end{cases}
\]
莫比乌斯函数 $\mu$、约数个数函数
\[
  \tau(n)=\sum_{d\mid n}1,
\]
约数和函数
\[
  \sigma(n)=\sum_{d\mid n}d,
\]
以及欧拉函数 $\varphi$。

狄利克雷卷积定义为
\[
  (f*g)(n)=\sum_{d\mid n}f(d)g\!\left(\frac nd\right).
\]
单位元是 $\varepsilon$，即 $f*\varepsilon=f$。
重要恒等式为
\[
  \mu*\mathbf 1=\varepsilon,
  \qquad
  \varphi*\mathbf 1=\operatorname{id},
  \qquad
  \mathbf 1*\mathbf 1=\tau,
  \qquad
  \operatorname{id}*\mathbf 1=\sigma.
\]

\begin{icpcresult}[莫比乌斯反演]
  若
  \[
    F(n)=\sum_{d\mid n}f(d),
  \]
  则
  \[
    f(n)=\sum_{d\mid n}\mu\!\left(\frac nd\right)F(d)
        =\sum_{d\mid n}\mu(d)F\!\left(\frac nd\right).
  \]
\end{icpcresult}

指示式
\[
  [\gcd(i,j)=1]=\sum_{d\mid i,\,d\mid j}\mu(d)
\]
常用于把互质限制变成可按倍数统计的整除条件。

\begin{icpcwarning}[卷积方向]
  反演公式有两种等价下标写法。
  使用时先明确外层函数是“对约数求和”还是“对倍数求和”，不要直接套用另一种形式。
\end{icpcwarning}
